\chapter{Conclusion}
\label{ch:conclusion}

This thesis presented a comprehensive framework for building production-ready LLM-powered natural language interfaces to spatial SQL databases, specifically addressing the accessibility challenge in urban energy modeling through the integrated CIM Wizard and CIM Assist systems.

\section{Contributions, Achievements, Impacts and Technical Insights}

This research makes substantial contributions across platform architecture, dataset generation methodology, fine-tuning efficiency, and evaluation frameworks, while demonstrating production-ready performance that validates the practical viability of LLM-powered spatial database interfaces. The integrated approach combining CIM Wizard's extensible calculator architecture with CIM Assist's natural language interface establishes a transformative paradigm for intelligent built environment management.

The \textbf{CIM Wizard Integrated platform} represents the first object-oriented framework specifically engineered for extensible urban energy modeling, fundamentally reimagining how researchers interact with complex spatial databases. Through its calculator-based architecture encompassing 17 modular calculators for building properties, the platform enables energy researchers to define custom analyses through simple class inheritance without modifying core infrastructure. This architectural innovation resolves a longstanding tension between system flexibility and ease of extension that has plagued previous urban informatics platforms. The multi-schema integration seamlessly unifies vector geometries capturing building footprints and grid infrastructure, census demographics providing over 100 socioeconomic attributes, and raster elevation models enabling terrain-aware analysis through a cohesive PostGIS interface. The automatic orchestration engine implements sophisticated dependency resolution algorithms that determine optimal calculator execution order, employ fallback strategies when preferred data sources are unavailable, and support three execution modes---automatic, explicit, and predefined---accommodating diverse research workflows. Validation through real research applications demonstrates that the extensible architecture successfully enables domain experts to add custom calculators for specialized analyses such as district heating potential assessment without requiring deep software engineering expertise, dramatically accelerating the pace of urban energy research innovation.

The \textbf{dataset generation pipeline} achieves unprecedented success through its two-stage approach that systematically builds quality and diversity while maintaining cost-effectiveness. Stage 1 employs 52 handcrafted templates across 11 SQL operation types with systematic parameter variation, generating 10,800 samples that establish a 98-100\% NoErr baseline through comprehensive schema validation addressing PostgreSQL-specific requirements including case-sensitive column references, explicit type casting for mathematical functions, correct ID column usage across different entity types, and precise project\_id scope awareness. Stage 2 applies multi-strategy LLM augmentation using GPT-4o-mini to generate diverse natural language questions through five complementary strategies (question paraphrase, SQL-to-question, question-to-SQL, SQL rewrite, and parameter variation), producing 50,000 augmented samples while achieving a remarkable \textbf{97.34\% NoErr rate} through schema-aware generation. This breakthrough validates that deep schema understanding through explicit constraints---including TABLE\_ID\_COLUMNS mappings preventing generic ID assumptions, TABLES\_WITH\_PROJECT\_ID whitelists restricting scope to only two relevant tables, strict complexity controls matching table counts to difficulty levels, and geometry awareness preventing invalid spatial operations on non-geometric tables---enables synthetic data generation to match template-based quality. The LLM augmentation stage completed in 39 hours at a total cost of merely \textbf{\$50}, achieving 99.7\% quality acceptance while demonstrating that task-appropriate model selection focusing expensive models only on truly difficult generation tasks yields substantial cost savings compared to human annotation.

The \textbf{validation-first methodology} represents a critical methodological contribution with implications extending far beyond this specific application domain. By validating quality at each stage before proceeding to the next rather than adopting the conventional sequential generate-then-validate approach, this work achieves high curation retention producing 34,993 training samples from 50,000 augmented samples. The mechanism underlying this success is elegant: Stage 2 samples achieve 97.34\% NoErr rate, meaning the challenging task of ensuring query executability is resolved during generation through schema-aware prompting. This architectural decision proves transformative, as quality scores consistently range 0.85-0.92 well above the 0.75 acceptance threshold, and samples marked with \texttt{no\_error: true} flags pass validation filters that would otherwise reject incorrectly generated alternatives. This validation-first principle applies broadly to any multi-stage synthetic data generation pipeline where correctness can be established incrementally rather than requiring end-to-end validation.

The \textbf{fine-tuning implementation} successfully demonstrates that parameter-efficient methods enable production-ready model adaptation on consumer hardware. Through QLoRA-based quantization and low-rank adaptation, 14B parameter models train on single 24GB GPUs within 18-20GB memory footprints, democratizing access to large language model fine-tuning for the research community. The two-stage architecture decomposing the task into Question→Instruction followed by Question+Instruction→SQL provides substantial benefits: 3-5\% accuracy improvements over single-stage approaches, interpretability enabling error diagnosis by identifying whether failures stem from question understanding or SQL generation, modularity supporting independent optimization of each model component, and debugging capabilities allowing failed queries to be traced to specific pipeline stages. Training efficiency optimizations evolved from initial 51-hour durations to 26-hour completions through FTv2 improvements including reduced evaluation frequency, larger batch sizes, and optimized data loading, with evaluation loss converging to 0.0972---a 2.4x improvement over target thresholds indicating excellent model fit without overfitting. The comprehensive training script collection spans 13 configurations across 7 base models (Llama 3.1, Qwen 2.5, DeepSeek, Mistral, Phi-3.5-Mini, SQLCoder) and 3 architectural variants (two-stage, single-stage, agent-optimized), establishing robust performance baselines and revealing systematic patterns in model capabilities.

Performance evaluation validates \textbf{production readiness} through systematic assessment across multiple complementary metrics. NoErr measuring dataset quality achieves 97.34\% for Stage 2 LLM augmentation, confirming that synthetic samples execute successfully against the actual database. Execution Accuracy (EX) assessing first-shot practical correctness demonstrates \textbf{84.58\% success rate} for SQLCoder 7B, substantially exceeding the 70\% production readiness threshold and representing 54 percentage points improvement over frontier models like GPT-4o that achieve only 30.85\% accuracy. Deep EM measuring structural correctness reaches 81.59\%, while Semantic Correctness (SC) achieves 63.18\%. Agent-based self-correction shows minimal improvement with EA matching EX at 84.58\%, indicating high first-shot quality where iterative refinement provides diminishing returns. The average iteration count of 1.11 confirms that 89\% of queries succeed on first attempt, with acceptable inference latency of 1.2 seconds remaining within practical bounds for interactive research applications.

The \textbf{technical insights} emerging from this work provide valuable guidance for future research in synthetic data generation and domain-specific model adaptation. The LLM-based augmentation achieving 97.34\% NoErr rate proves that large language models can reliably produce complex structured data when properly constrained by explicit schema knowledge through carefully designed prompts. The task-appropriate model selection strategy employing cost-effective GPT-4o-mini for question generation achieves order-of-magnitude cost reductions at \$50 for 50,000 samples, demonstrating that careful task decomposition enables practical large-scale dataset creation. The validation-first methodology transforms quality assurance through stage-wise validation without algorithmic innovation, highlighting that methodological decisions often matter as much as technical sophistication. The direct Q2SQL fine-tuning approach provides optimal balance between accuracy (84.58\%) and inference speed (1.2 seconds), validating that specialized pre-training (SQLCoder on BIRD benchmark) combined with domain-specific fine-tuning outperforms larger general-purpose models. Agent integration with database introspection tools reveals that fine-tuned models achieve high first-shot quality with minimal self-correction benefit, validating that quality training data proves more valuable than iterative refinement mechanisms.

The \textbf{broader impact} of this framework extends well beyond urban energy modeling to fundamentally reshape how researchers and practitioners interact with complex spatial databases across diverse domains. This work establishes that intelligent built environments---whether focused on urban energy systems, transportation networks, environmental monitoring, or infrastructure planning---require sophisticated information model generators that maintain structured representations of spatial objects with their attributes, relationships, and computational methods. CIM Wizard provides precisely this foundation through its calculator-based architecture, enabling buildings, grid infrastructure, census regions, and environmental sensors to exist as first-class objects with well-defined properties computed through multiple fallback methods ensuring robustness. CIM Assist then layers natural language accessibility atop this structured foundation, transforming the platform from an expert-oriented database system into an intelligent environment where domain specialists query complex spatial relationships through conversational interfaces. This integration proves \textbf{essential for any truly intelligent built environment}, as neither component suffices independently: structured object models without accessible interfaces remain locked behind technical barriers, while natural language interfaces without robust underlying data models lack the semantic richness necessary for sophisticated analysis.

The democratization of spatial analytics enabled by this work reduces barriers for urban planners analyzing building energy without SQL expertise, energy consultants rapidly prototyping analysis workflows, policymakers making data-driven decisions without technical dependencies, and researchers focusing on domain questions rather than database queries. The \$50 cost for 50,000 augmented samples compared to \$50,000-100,000 for equivalent human annotation demonstrates that high-quality domain-specific datasets can be generated at scale without prohibitive costs, making advanced AI techniques accessible to research communities worldwide. The complete open-source implementation published under Apache 2.0 license---encompassing CIM Wizard platform, AI4DB dataset generation pipeline, TXT2SSQL fine-tuning scripts, CIM Assist agent integration, and the HuggingFace taherdoust/ai4cimdb dataset---enables reproducibility, extensibility to new domains, community collaboration, and serves as an educational resource for Text-to-SQL and spatial database courses.

This work advances Text-to-Spatial-SQL research by providing the first comprehensive framework with production-ready performance exceeding 80\% accuracy thresholds, establishing evaluation benchmarks and datasets for future comparisons, setting performance baselines through fine-tuned open-source models, quantifying the 20-30\% PostGIS knowledge gap in generic models, and creating a foundation for multi-domain spatial SQL research spanning transportation, environment, healthcare, and urban planning applications.

\section{Limitations and Future Directions}

While this research demonstrates production-ready performance within its intended scope, several limitations suggest valuable directions for future investigation that could substantially enhance both the CIM Wizard platform capabilities and the CIM Assist interface generalizability, efficiency, and applicability across broader contexts.

\subsection{Model Performance Limitations}

Despite strong performance achievements, the evaluation reveals specific areas requiring improvement. The 84.58\% execution accuracy, while exceeding production readiness thresholds, leaves room for enhancement particularly on complex multi-schema queries where accuracy drops to 72\%. These challenging queries involving joins across cim\_vector, cim\_census, and cim\_raster schemas with spatial predicates represent the frontier of model capability, suggesting that specialized training on multi-schema patterns or architectural innovations could yield substantial gains.

The evaluation benchmark of 201 samples, though carefully stratified across task types and complexity levels, may not capture all production edge cases that emerge during real-world deployment. Continuous evaluation as usage patterns evolve will prove essential for maintaining performance guarantees, particularly as users explore novel query patterns not anticipated during initial dataset generation. Future work should explore ensemble approaches combining fine-tuned models with frontier models' reasoning capabilities, potentially achieving higher accuracy on complex queries through complementary strengths. Integration of retrieval-augmented generation (RAG) could provide dynamic schema updates without retraining, addressing the challenge of evolving database structures common in research environments. Investigation of few-shot prompting techniques might enable rapid adaptation to schema changes while maintaining base model performance, reducing the need for complete retraining cycles when new tables or columns are added.

\subsection{CIM Wizard Platform Extensions}

The CIM Wizard platform, while successfully demonstrating the calculator-based architecture's viability for urban energy modeling, requires substantial extensions to achieve comprehensive intelligent built environment capabilities. The current implementation operates within a relatively constrained spatial representation that limits its applicability to advanced building energy analysis workflows employed by the broader research community. Integration with standardized urban information models, particularly 3DCityDB schema implementing CityGML standards, represents a critical next step enabling interoperability with the extensive ecosystem of urban modeling tools and data sources that have adopted these international standards. CityGML's multi-level-of-detail paradigm spanning from coarse city-scale representations (LOD0-LOD1) through detailed architectural models (LOD3-LOD4) would enable the platform to serve diverse analysis scales simultaneously, supporting both city-wide policy analysis and building-specific energy optimization within a unified framework. This deeper spatial granularity incorporating building interior geometries, thermal zones, construction assemblies, and envelope characteristics would dramatically expand analytical capabilities while introducing schema complexity that reinforces the necessity of natural language interfaces like CIM Assist for maintaining user accessibility.

The calculator library currently encompassing 17 building property methods requires expansion to align with established building energy analysis standards, particularly the TABULA building typology framework (https://webtool.building-typology.eu/) that provides standardized residential building archetypes across European countries. Implementing TABULA-compatible calculators would enable seamless integration with existing building stock analysis methodologies, facilitate cross-country comparative studies, and ensure consistency with policy-relevant energy assessment frameworks employed by European energy agencies. Additional calculators addressing heating system efficiency, ventilation strategies, domestic hot water consumption patterns, occupant behavior modeling, and retrofit measure impact assessment would position CIM Wizard as a comprehensive platform supporting the full building energy analysis workflow from initial assessment through detailed optimization and policy evaluation.

Perhaps most significantly, the current database architecture lacks temporal data management capabilities essential for storing and analyzing time-series energy simulation results that form the core output of building performance simulation tools. Extending the PostGIS database with TimescaleDB—a PostgreSQL extension specifically optimized for time-series data—would enable CIM Wizard to host high-resolution temporal simulation outputs including hourly heating loads, minute-level power consumption profiles, seasonal energy patterns, and multi-year performance trajectories. This temporal extension transforms the platform from a static spatial database into a full tempo-spatial relational database capable of supporting sophisticated queries combining spatial relationships with temporal patterns, such as identifying building clusters exhibiting similar diurnal energy profiles or analyzing how building geometry influences seasonal heating demand variations. The resulting infrastructure would provide researchers with unprecedented analytical capabilities for understanding urban energy dynamics across both spatial and temporal dimensions within a unified database framework, though this enhanced complexity further emphasizes the critical importance of CIM Assist's natural language interface for maintaining analytical accessibility.

Critically, as CIM Wizard's schema expands to incorporate the deeper spatial levels of detail inherent in CityGML LOD3-LOD4 representations and temporal dimensions through TimescaleDB integration, the resulting database complexity will reach levels where CIM Assist transitions from a convenience feature to an absolute necessity even for domain-specific applications within urban energy modeling. The current schema with approximately 20 tables across three schemas already challenges non-technical users; expanding to hundreds of tables representing building interior geometries, construction assemblies, thermal zones, HVAC systems, and time-series simulation results at multiple temporal resolutions would create a database structure of ultra-high complexity that even SQL-proficient domain experts would struggle to navigate efficiently. At this complexity threshold, formulating correct queries requires simultaneously understanding spatial hierarchies, temporal relationships, foreign key dependencies, and domain-specific calculation methods—a cognitive load that fundamentally exceeds human working memory capacity without intelligent assistance. CIM Assist's natural language interface, enhanced through continuous learning from expert corrections, would become the primary access modality enabling researchers to leverage the platform's full analytical capabilities without drowning in schema complexity. This evolution toward schema-complexity-driven interface necessity represents a general pattern for intelligent built environments: as information models grow richer and more semantically detailed to support sophisticated analysis, natural language interfaces transition from optional conveniences to essential infrastructure components.

\subsection{CIM Assist Interface Limitations}

The CIM Assist natural language interface, while demonstrating production-ready performance on the current CIM Wizard schema, faces generalization challenges and infrastructure limitations that constrain its broader applicability. The current evaluation focuses exclusively on the CIM Wizard database, establishing strong performance in urban energy modeling contexts while leaving generalization to other spatial databases empirically unvalidated. Performance on transportation networks analyzing road infrastructure and public transit, environmental monitoring systems tracking pollution and climate data, or healthcare facility databases assessing accessibility and service coverage remains unconfirmed, though the domain-agnostic nature of PostGIS operations and the modular calculator architecture suggest reasonable transferability potential. The models learn CIM-specific terminology including building types, energy metrics, and urban infrastructure concepts that may not directly transfer to other domains, and the spatial function coverage focuses on operations prevalent in urban energy analysis potentially underrepresenting functions critical to domains like network routing in transportation or hotspot analysis in epidemiology. Future multi-database evaluation should systematically assess the framework across diverse spatial domains including transportation infrastructure planning with routing and network analysis operations, environmental resource management with buffer zones and temporal trend analysis, healthcare facility accessibility with service areas and proximity analysis, and urban planning with zoning overlays and suitability assessment. This evaluation would validate generalization capabilities, identify domain-specific patterns requiring specialized treatment, and develop cross-domain transfer learning strategies that leverage common PostGIS knowledge while adapting to domain-specific vocabularies.

The Stage 2 dataset generation bottleneck requiring 39 hours for completion through sequential API calls at 2.25 seconds per sample creates a practical limitation restricting rapid iteration during dataset refinement cycles. This linear time dependency scales poorly with dataset size increases, limiting experimentation with larger training corpora or alternative augmentation strategies. The conservative sequential architecture implemented to simplify error handling and checkpoint management during initial development sacrifices throughput for reliability, representing a reasonable engineering trade-off for research development but suboptimal for production deployment. Implementing parallel batch processing with intelligent rate limiting could reduce total generation time to 30-60 minutes by exploiting modern API infrastructure designed for concurrent request handling, representing approximately 100x speedup that would transform multi-hour waits into rapid iteration cycles. The technical approach would involve batching 100-200 requests per minute respecting API rate limits, implementing adaptive rate limiting based on API response headers and error codes, employing asyncio for concurrent processing across multiple connections, and developing retry logic for gracefully handling rate limit errors. This optimization faces the practical challenge of balancing speed gains against OpenRouter rate limiting policies to avoid throttling or service disruptions, requiring careful batch sizing strategies and request pacing algorithms that maximize throughput while maintaining service reliability.

Training time overhead, while substantially reduced from initial 51-hour durations to 26-hour completions through FTv2 optimizations, still limits rapid experimentation and hyperparameter exploration critical for research advancement. The 43.6\% time spent on evaluation during training---165 total evaluation runs checking validation loss and model quality---represents a necessary investment for monitoring convergence but consumes substantial computational resources. Further optimizations targeting 15-20 hour training durations for 14B models appear achievable through gradient checkpointing providing 15-20\% memory reduction enabling larger batch sizes, Flash Attention 2 accelerating attention computation with 20-30\% speedup, mixed precision training leveraging FP16/BF16 for 15-25\% speedup, and larger batch sizes exploiting improved GPU utilization for 10-15\% additional speedup. These optimizations would enable faster experimentation cycles, more extensive hyperparameter searches, and evaluation of additional model variants without proportionally increasing computational budgets.

Hardware constraints limiting experimentation to single 24GB GPU configurations restrict exploration of larger model architectures that might provide additional accuracy gains. While 14B parameter models achieve production-ready performance, evaluating 32B models through Unsloth framework optimizations or 70B models through multi-GPU distributed training could reveal whether larger architectures provide sufficient 5-10\% accuracy improvements to justify 2-3x increases in training time and inference latency. The trade-off analysis must carefully weigh accuracy gains against computational costs and deployment complexity, with recommendations varying based on whether accuracy priority outweighs efficiency considerations in specific application contexts.

\subsection{Extended Future Work Directions}

Implementing continuous learning mechanisms through user feedback loops would transform CIM Assist from a static trained model into an adaptive system that improves through deployment experience. Capturing user corrections to generated SQL queries provides high-quality training signals that directly reflect real-world usage patterns and error modes not fully represented in synthetic training data. This feedback can drive a semi-reinforcement learning pipeline where the agent receives reward signals based on whether database execution succeeds and whether users accept or modify generated queries, enabling incremental model improvement without requiring expensive manual annotation. Active learning strategies would prioritize uncertain queries where the model exhibits low confidence or generates multiple competing interpretations, focusing human feedback on cases that provide maximum informational value for model refinement. This continuous adaptation mechanism becomes particularly valuable as the CIM Wizard schema evolves, enabling the natural language interface to rapidly accommodate new tables, columns, and calculators through few-shot learning on user-corrected examples rather than requiring complete model retraining cycles.

Future research directions extending beyond the immediate limitations include cross-lingual extensions supporting Italian, Spanish, German, and French to broaden accessibility for European Union projects where multilingual support proves essential for stakeholder engagement. Multi-modal integration incorporating maps for point-and-click spatial region specification, building images for visual similarity queries through vision models like CLIP and LLaVA, and urban planning diagrams for diagram-based query construction would dramatically enhance user experience. Federated learning approaches enabling collaborative model training across institutions without data sharing would address privacy concerns prevalent in urban informatics while improving generalization through exposure to diverse data sources. These extensions would further solidify CIM Assist's role as the essential human-machine interface for next-generation intelligent built environment platforms.

The modular architecture, comprehensive evaluation framework, and substantial open-source implementation provide a solid foundation for these future research directions. The validation-first methodology, schema-aware LLM augmentation principles, and task-appropriate model selection strategies established through this work offer guidance applicable well beyond the specific urban energy modeling domain, suggesting that the contributions extend to the broader challenges of building reliable, cost-effective, production-ready natural language interfaces for complex structured databases across diverse application domains.

